\subsection{Retaining all prime-field components}
A homogeneous linearized relation with zero coefficient sum does not
specify one prime-field value of its argument. It nevertheless permits
a smaller equation that keeps every such value. Additive composition
is classical \citep[\S6]{sp-giesbrecht}; the issue here is recognizing
the hidden argument and controlling support through repeated reductions.

\begin{theorem}[A zero-sum relation with controlled support growth]
\label{sp:thm:zero-sum-preprocessing}
Fix $p$ and let a homogeneous sparse binary-exponent list generate
$K\subset\Fp[x_0,\ldots,x_N]$, with $I=K^{\mathrm{sat}}$ and $S/I$
geometrically integral. Suppose a nonzero row admits
\[
 F=U^Q+\sum_{i=0}^{s}a_i V^{Q-p^i}U^{p^i},
 \qquad Q=p^e>p,\quad 0\le s<e,
\]
where $U$ is homogeneous, $V$ is a coefficient-one monomial of the
same degree $d\ge1$, and
\[
 (d+1)p^s\le Q,\qquad a_i\in\Fp,
 \qquad 1+\sum_i a_i=0.
\]
None of these data is given. A deterministic polynomial sparse-bit
algorithm recognizes and exactly verifies such a representation whenever
one exists, normalizes $U$ to exclude the monomial $V$, and replaces
$F$ by the nonzero polynomial
\begin{equation}
 G=U^p-V^{p-1}U.
 \label{sp:eq:zero-sum-output}
\end{equation}
It satisfies
\[
 K\subseteq K'\subseteq I,
 \qquad (K')^{\mathrm{sat}}=I,
 \qquad \deg G=\frac{p\deg F}{Q}<\deg F.
\]
No coordinate index is added, and the output has at most twice the
input support. A deterministic convention can minimize the output
degree among passing representations. Despite this possible increase,
repeated use has polynomial total sparse-bit cost, also when
interleaved with the preceding coordinate, digit, linearized and
affine reductions. If $T_0$ is the original total support and $D_0$
the largest original degree, the total support throughout any such
interleaved computation is at most
\[
 16T_0\max\{1,\log_p D_0\}.
\]
\end{theorem}
\begin{proof}
Put $D=\deg F$. The translation $U\mapsto U-cV$, $c\in\Fp$,
leaves the relation unchanged because its coefficient sum is zero.
Remove the $V$-term of $U$. The normalized $U$ is nonzero, since
otherwise $F=0$. The same gap calculation as before gives
\[
 U=E_Q(F),
\]
where $E_Q$ is the divided coordinatewise $Q$-section. Indeed,
$Qv+p^i(u-v)$ is coordinatewise divisible by $Q$ only when $u=v$.
Thus the inner polynomial is determined at every candidate scale.

Enumerate powers $Q>p$ dividing $D$, and the admissible values of
$s$, all in polynomial number in $\log(D+1)$. Set $d=D/Q$, reject a
zero section, and form $R=F-U^Q$. A genuine normalized representation
has $R\ne0$. Otherwise $U/V$ in the rational function field would
satisfy the nonzero polynomial $\sum_i a_iT^{p^i}$, because
$\sum_i a_i=-1$. The polynomial ring is geometrically integral, so
its algebraic constants are $\Fp$. This would make $U$ a scalar
multiple of $V$, contrary to its normalization.

For completeness of the monomial search, put $q=p^s$ and $M=Q-q$,
and choose any exponent $\gamma$ in $R$. At least one contributing
lower-block term has
\[
 \gamma=(Q-p^i)v+p^iu
        =Mv+\beta,
 \qquad \beta=(q-p^i)v+p^iu\ge0,
 \qquad |\beta|=qd\le M.
\]
Therefore the coordinatewise floor $\lfloor\gamma/M\rfloor$ equals
$v$, or equals $v+e_j$ for a single coordinate $j$. Enumerate these
possibilities and keep the nonnegative vectors of degree $d$ whose
monomial is absent from $U$. This includes the equality case $qd=M$.
For each candidate solve
\[
 \sum_{i=0}^{s}a_iV^{Q-p^i}U^{p^i}=R,
 \qquad \sum_{i=0}^{s}a_i=-1
\]
by coefficientwise linear algebra, and verify the exact sparse identity.
There are $O(\log D)$ unknowns and
$O(|\operatorname{supp}F|\log D)$ equations. Frobenius powers and
monomial shifts construct every column without dense expansion.
The true representation passes. All accepted representations meet
the stated conditions. Decreasing order of $Q$ minimizes
$pD/Q$, with fixed further orders resolving ties. Every coordinate
of $V$ occurs in $R$, since its lower-block exponents are at least
$Mv$ coordinatewise; hence no new index is introduced.

Let $\overline U,\overline V$ be the images in $S/I$. If
$\overline V=0$, the original relation gives $\overline U=0$.
Otherwise $\overline U/\overline V$ is algebraic over $\Fp$,
hence belongs to $\Fp$. In either case $G\in I$.
To prove the other ideal inclusion, write
\[
 L(T)=\sum_{k=0}^e c_kT^{p^k},\qquad
 c_e=1,\quad \sum_{k=0}^e c_k=0,
 \qquad b_j=-\sum_{k=0}^j c_k.
\]
Here $c_k=a_k$ for $k\le s$ and is zero at the other lower indices.
Adjacent coefficients give
\[
 L(T)=\sum_{j=0}^{e-1}b_j(T^p-T)^{p^j}.
\]
Homogenizing to degree $Q$ in $U,V$ yields the exact identity
\begin{equation}
 F=\sum_{j=0}^{e-1}b_j
       V^{Q-p^{j+1}}G^{p^j}.
 \label{sp:eq:zero-sum-divisibility}
\end{equation}
All exponents of $V$ are nonnegative. Thus $F\in(G)$, proving
$K\subseteq K'\subseteq I$ and equality after saturation.
The polynomial $G$ is nonzero: if it vanished in the ambient
polynomial ring, the nonconstant rational function $U/V$ would be a
root of $T^p-T$, again impossible. Formula~\eqref{sp:eq:zero-sum-output}
gives its degree and its support bound.

It remains to control iteration. Let $T_0$ be the original total
support, $D_0$ the largest original degree, and $v$ the number of
initially occurring coordinates. Every current nonzero row belongs
to the proper homogeneous ideal $I$, so its degree is positive.
The gap and $Q>p$ give
\[
 D=dQ>d^2,\qquad D_{\mathrm{new}}=pd<p\sqrt D,
 \qquad D_{\mathrm{new}}\le D/p.
\]
Along a lineage put $z=\log_p D$. At a zero-sum step,
\[
 z_{\mathrm{new}}<1+\tfrac12z_{\mathrm{old}},
 \qquad z_{\mathrm{new}}\le z_{\mathrm{old}}-1.
\]
Intervening earlier operations only decrease degree. After $j\ge1$
zero-sum steps the first inequality therefore gives
$z_j<2+(z_0-2)/2^j$. After
$\lceil\log_2\max\{1,z_0\}\rceil$ such steps, if that many occur,
the degree is below $p^3$; at most three further steps are possible
by the second inequality. Thus every lineage has at most
\[
 k_0=\left\lceil\log_2\max\{1,\log_p D_0\}\right\rceil+3
 \quad\text{zero-sum steps},\qquad
 2^{k_0}\le16\max\{1,\log_p D_0\}.
\]

Give each original term a formal token. Across a
support-nonincreasing operation, assign all new terms injectively
to old tokens, including terms in different affine output rows.
Across a zero-sum step assign at most two new tokens to each old
token. Unassigned tokens terminate. These assignments only count
terms and assert no matching of polynomial coefficients. A token
follows the degree of its current row. Branching occurs only at a
zero-sum step, so each initial token has at most $2^{k_0}$ descendant
leaves, including terminated tokens. The current tokens form an
antichain in this forest. Consequently the total support, across
all rows and all splits, satisfies
\begin{equation}
 T_{\max}\le T_0 2^{k_0}
 \le16T_0\max\{1,\log_p D_0\}.
 \label{sp:eq:zero-sum-support-budget}
\end{equation}

To count all non-coordinate operations, suppress only
coordinate-content transitions in this forest and retain every
other transition, including unary transitions of earlier rules.
Every nontrivial remaining step makes each output degree at most
$\rho=1-1/(v+2)$ times its input degree: this is the digit guarantee,
while the other steps divide degree by at least $p$. No operation
introduces a coordinate, so the active count never exceeds $v$.
Since $-\log\rho\ge1/(v+2)$, every root-to-leaf path has at most
\[
 h_0=\left\lfloor(v+2)\log D_0\right\rfloor
\]
edges. A tree with at most $2^{k_0}$ leaves and height $h_0$ has at
most $2^{k_0}(h_0+1)$ vertices, because its root-to-leaf paths cover
every vertex. Each successful non-coordinate row operation has a
nonzero output and can be charged to one of its new token edges;
distinct operations create distinct edges. The total number of
these operations is therefore
\[
 O\!\left(T_0(v+2)(1+\log D_0)^2\right).
\]
This height-times-leaves count includes all unary earlier-rule
transitions, so no vertex bound is inferred from branching depth
alone.

Perform maximal eligible coordinate-content cleanup initially and
between successive non-coordinate operations. Removing one
coordinate's complete content does not create content in another
coordinate, so a batch has at most $v$ removals per current row.
The number of rows is at most $T_{\max}$, and the number of batches
is bounded by one plus the operation count just proved. Thus
coordinate cleanup also has polynomial cost. Coefficients remain
in $\Fp$, coordinate indices never grow, and exponent lengths are
bounded by the original degree bit length. Every intermediate list
has polynomial sparse-bit size. Complete recognition scans have
polynomial cost in that size; a successful scan is charged to a
successful operation or cleanup batch, and a complete unsuccessful
scan terminates. This proves the full interleaved bit bound.
\end{proof}

Support can increase in one allowed step. Over $\F_3$, let $m\ge3$,
$d=3^{2(m-1)}$, $V=y^d$, and
\[
 U=\sum_{j=0}^{m-1}(-1)^j
       x^{3^{2j}}y^{d-3^{2j}},\qquad Q=3^{2m+1}.
\]
Take $a_0=a_2=1$, with the other lower coefficients zero. The two
lower blocks telescope to two endpoint terms, so $F$ has $m+2$
terms. The two exponent sets in $U^3-V^2U$ consist of odd and even
powers of three and are disjoint, giving $2m$ terms. At $m=3$ this
is a five-term input and a six-term output. A different passing
representation may give a smaller output; the example shows why a
uniform iteration proof cannot simply assume support nonincrease.

\begin{example}[A strict family for the zero-sum reduction]
\label{sp:ex:zero-sum-family}
Over $\F_2$, let $n\ge5$ be odd and put
\[
 g=\sum_{j=0}^{n-1}x_j^n+x_0^{n-1}x_1,
 \qquad \mathcal M=\operatorname{supp}(g),
 \qquad Q=2^r>n.
\]
Here a member of $\mathcal M$ denotes its coefficient-one monomial.
Let $K$ be generated by the displayed rows
\begin{equation}
 g^Q-V^{Q-1}g,\qquad
 g^{2Q}-V^{2Q-1}g,
 \qquad V\in\mathcal M.
 \label{sp:eq:zero-sum-family}
\end{equation}
Then $K^{\mathrm{sat}}=(g)$, whose projective hypersurface in
$\PP^{n-1}$ is smooth, geometrically integral and anticanonical.
The preceding coordinate, digit, nonzero-sum and affine reductions
leave this displayed list unchanged, and its least principal
complete-intersection excess is $n(Q-1)$.
Theorem~\ref{sp:thm:zero-sum-preprocessing}, choosing minimum output
degree, replaces its rows by $g(g+V)$, of degree $2n$ and with an
$n$-term degree-$n$ cofactor. The designated sparse-cofactor algorithm
therefore computes the scalar in polynomial bit time.
For $n=2r+1$, $r\ge3$, the input length is polynomial in $r$ while
the earlier residual numerical cofactor degrees are exponential in $r$.
\end{example}
\begin{proof}
The partial derivatives of $g$ are
\[
 g_{x_0}=x_0^{n-1},\qquad
 g_{x_1}=x_1^{n-1}+x_0^{n-1},\qquad
 g_{x_j}=x_j^{n-1}\quad(j\ge2).
\]
Their only common affine zero is the origin, proving projective
smoothness. Positive-degree factors would intersect projectively
and yield a singular point; repeated factors are likewise excluded.
Thus $g$ is geometrically irreducible. Its degree is the number $n$
of homogeneous coordinates, so it is anticanonical.

The cofactors $g^{P-1}-V^{P-1}$, $P=Q,2Q$, have no common
projective zero. Where $g=0$, vanishing of all cofactors forces all
$V$ to vanish, including every $x_j^n$. Where $g\ne0$, their
vanishing gives
\[
 (V/g)^{Q-1}=(V/g)^{2Q-1}=1.
\]
The two exponents are coprime, so $V/g=1$ for all $V\in\mathcal M$.
But $|\mathcal M|=n+1$ is even, contradicting
$1=g/g=\sum_{V\in\mathcal M}V/g$. The cofactor ideal sheaf is
therefore the unit sheaf, proving the scheme-theoretic saturation.
The least nonzero degree in $K$ is $nQ$, attained by a row of
scale $Q$, so its principal excess is $nQ-n$.

Fix a row of scale $P=Q$ or $2Q$. Its pure top terms include
$x_j^{nP}$ for every pure monomial different from $V$.
Only $V^P$ cancels between the two blocks: any further cancellation
would require $P(\alpha-v)=u-v$, with every coordinate on the
right of absolute value less than $P$. At least $n-1$ pure top
terms remain, so the monomial content is one.

Consider a digit scale $A=2^a\le P$. Every pure top term has low
digit zero. Choose a pure monomial $x_j^n$ whose coordinate is
outside the support of $V$. Its mixed term has $j$-exponent $n$,
which is odd, and therefore gives a nonzero low row. Since the
common total degree $nP$ is divisible by $A$, a column of high
total degree $nP/A$ can occur only in low row zero. A surviving
pure top column and any nonzero entry in the nonzero low row form
a triangular two-by-two minor with nonzero diagonal. The digit
rank is at least two.

If $P<A\le nP$, write $A=Pb$, with $b$ a power of two satisfying
$2\le b\le n$. Choose two indices $j\ge2$ whose pure monomials
are not $V$; $n\ge5$ guarantees two choices. Their top terms have
distinct low rows $P(n\bmod b)e_j$ and distinct high columns
$\lfloor n/b\rfloor e_j$. Both factors are nonzero because $n$
is odd and $b\le n$. They form an identity minor. The cross
monomial involves only indices zero and one. Every remaining mixed
term has a coordinate not divisible by $P$, since $P>n$ and
$u\ne v$; it cannot enter those low rows. Thus the minor persists.
For $A>nP$ the high factor is constant. No nontrivial digit
reduction applies at any scale.

In characteristic two a nonzero additive coefficient sum equals
one. Its top-section formula forces the inner polynomial of any
passing earlier relation at scale $A$ to be $E_A(F)$. Because $n$
is odd and $A$ divides $nP$, necessarily $A\le P$. Here
\[
 E_A(F)=g^{P/A}-H,
 \qquad H=E_A(V^{P-1}g).
\]
The term $V$ contributes the nonzero monomial $V^{P/A}$ to $H$.
The other pure monomials have different residues modulo $A$.
The cross monomial shares a residue with a pure monomial only for
$x_1^n$, and then only if $A\mid n-1$. Thus $H$ is nonzero and
involves either one coordinate or only $x_0,x_1$. It cannot be a
multiple of $g$, which has positive degree in every other coordinate.
Consequently $E_A(F)\notin(g)$, contradicting the safety conclusion
of any passing nonzero-sum relation. An affine no-root relation is
also impossible: it would put a positive-degree monomial in $(g)$,
although this prime ideal contains no coordinate. None of the
earlier rules changes the row, so their whole displayed-row closure
leaves the list unchanged.

For the new rule take $U=g+V$, $s=0$, and $a_0=1$. The degree gap
is $n+1\le P$, and its output is
\[
 (g+V)^2-V(g+V)=g(g+V).
\]
No smaller output degree is possible, since $P$ is the largest
power of two dividing $nP$. At this scale $U=E_P(F)=g+V$ is fixed.
If a different monomial $W$ gave a passing output, safety modulo
$(g)$ would imply $V(V+W)=0$, hence $V+W\in(g)$. This is a
polynomial of degree $n$ with at most two terms, whereas $g$ has
$n+1\ge6$ terms. It must be zero, so $W=V$. Thus the
minimum-degree output is precisely the stated row. Its cofactor
$g+V$ has degree $n$ and $n$ terms, meeting
Theorem~\ref{sp:thm:sparse-cofactor} with explicit polynomial bounds.

Finally, there are $2(n+1)$ original rows, each with $2n$ terms and
exponent bit length $O(r+\log n)$. Setting $n=2r+1$ gives polynomial
input length in $r$, while $n(Q-1)$ is exponential. This comparison
concerns the specified preprocessing and cofactor conditions; it
asserts no hardness of the scalar on this family.
\end{proof}
